\documentclass[11pt,fontset=none]{ctexart}
\usepackage[margin=1in]{geometry}
\usepackage{booktabs}
\usepackage{enumitem}
\usepackage{xcolor}
\usepackage{hyperref}
\setCJKmainfont{Songti SC}
\setCJKsansfont{PingFang SC}
\setCJKmonofont{PingFang SC}

\hypersetup{
  colorlinks=true,
  linkcolor=blue,
  urlcolor=blue,
  citecolor=blue
}

\setlist[itemize]{leftmargin=1.7em,itemsep=0.25em,topsep=0.35em}
\setlist[enumerate]{leftmargin=1.7em,itemsep=0.25em,topsep=0.35em}

\title{RetailBench ARR Revision 与 paper/latex 版本差异说明}
\author{RetailBench Revision Notes}
\date{2026-05-25}

\begin{document}
\maketitle

\section{目的}

本文档说明当前 \texttt{paper/latex\_arr\_revision} 版本相对于
\texttt{paper/latex} 版本的主要差异。比较范围以两个目录中当前的 LaTeX
源码、表格、图表和参考文献文件为准，不包含外部实验原始日志。

\section{总体结论}

\texttt{paper/latex\_arr\_revision} 已经从早期的 ``simulator + method''
叙述转向 ``benchmark + evaluation + diagnostic analysis'' 叙述。新版主线更集中：
它强调 RetailBench 是一个 data-grounded、long-horizon、partially observable
retail operation benchmark，并围绕 environment construction、evaluation
protocol、main results 和 four-stage diagnostic analysis 展开。旧版
\texttt{paper/latex} 保留了更多早期章节结构，例如 simulator execution logic、
standalone results section、prompt/framework appendix，以及部分旧实验叙述。

\section{结构差异}

\begin{table}[h]
\centering
\small
\begin{tabular}{p{0.28\linewidth}p{0.31\linewidth}p{0.31\linewidth}}
\toprule
维度 & \texttt{paper/latex} & \texttt{latex\_arr\_revision} \\
\midrule
正文组织 & Introduction, Environment, Simulator Module Design, Experiment Settings, Results, Analysis, Related Work & Introduction, Environment Construction, Evaluation Protocol, Experiments, Analysis, Related Work \\
环境描述 & 更偏实现流程和模块执行逻辑 & 以 POMDP、Figure 1/Figure 2、数据源和模块设计为主线 \\
实验描述 & 分散在 Experiment Settings、Results、Analysis 中 & 合并为 evaluation protocol、main experiment results 和 four-stage diagnostics \\
附录 & 包含环境、rollout、prompt、agent framework 等多个 appendix & 当前主 PDF 保留环境配置与 state/action/module/update 细节，减少正文负担 \\
约束文件 & 无专门约束文件 & 新增 \texttt{paper\_constraint.md}，约束页数、Figure 位置、术语和 product 命名 \\
\bottomrule
\end{tabular}
\end{table}

\section{内容差异}

\begin{enumerate}
  \item \textbf{Benchmark positioning 更明确。}
  新版 introduction 和 related work 将 RetailBench 定位为 long-horizon retail
  operation benchmark，减少对短任务或通用 agent benchmark 的泛泛比较，并补充
  LifelongAgentBench 和 TheAgentCompany 等 long-horizon/workplace agent benchmark。

  \item \textbf{Environment construction 更结构化。}
  新版将环境描述拆为 benchmark、POMDP、data、modules 四个逻辑层次。正文只保留
  高层说明，详细 state space、action space、module definitions、day-level update
  equations 和 hand-crafted policy 放到 appendix。

  \item \textbf{Figure 设计被重排。}
  新版将 retail simulator module figure 作为 Figure 1 放在第一页右栏顶部，并将
  agent--simulator interaction framework 作为 Figure 2 双栏展示。旧版图表更接近初始
  draw.io 版本，缺少 shelf assortment 和当前交互框架的统一表达。

  \item \textbf{实验结果改为当前 paper-submit 数据。}
  新版实验部分基于 \texttt{paper\_submit\_data/outputs/report.md} 和四阶段分析报告，
  汇总 7 个模型、3 个 frameworks 以及 hand-crafted policy 的最佳结果。Metrics
  被重新组织为 viability、sales capability、inventory control、tool use 和 inference cost。

  \item \textbf{Analysis 改为 four-stage diagnostics。}
  新版 Analysis 围绕 evidence acquisition、evidence-to-action conversion、action quality
  和 long-horizon adaptation 四个 stage 展开，并配套单栏图表。旧版 analysis 更接近早期
  failure case 和行为观察汇总。

  \item \textbf{术语和表格做了收敛。}
  新版统一使用 frameworks，避免 scaffold；全局用 product/products 替代 SKU/SKUs；
  Table 1 被移入 appendix，并调整 benchmark 引用和列设计。
\end{enumerate}

\section{引用审查结果}

对当前 \texttt{latex\_arr\_revision} 的 \texttt{.tex} 文件进行引用抽取后，所有正文
citation key 都能在 \texttt{custom.bib} 中找到，未发现缺失 key。进一步检查 arXiv ID、
DOI 和数据集 URL 后，当前版本未发现明确的虚假引用。已做的保守修正包括：

\begin{itemize}
  \item 修正 McFadden 条目的 booktitle 拼写：\texttt{Fontiers} 改为 \texttt{Frontiers}。
  \item 为 Amazon Reviews 2023 对应的 \texttt{hou2024bridging} 条目补充 arXiv metadata 和 URL。
  \item MNL 模型引用中去掉弱相关的书评类引用 \texttt{Uncles01041987}，保留经典 McFadden 引用。
  \item 移除 revision 目录中未被主文档引用的旧 CAEC/FCR 文件，避免源码包内部出现过期方法叙述。
\end{itemize}

\section{提交建议}

建议提交 \texttt{paper/latex\_arr\_revision} 的 LaTeX source、最终 PDF、必要图表、
表格、参考文献和本文档 PDF；同时提交 \texttt{paper\_submit\_data} 中用于支撑当前实验结果的
分析脚本和输出报告。LaTeX 中间文件、\texttt{.DS\_Store}、\texttt{*.aux}、\texttt{*.log}、
\texttt{*.synctex.gz} 等构建产物不应提交。

\end{document}
